\section{Stochastic computing（随机计算）}

by B. R. GAINES

Standard Telecommunication Laboratories

Harlow, Essex, United Kingdom

\section{引言}

随机计算机是作为学习机形式的先进自动控制器之结构、实现与应用研究计划的一部分而发展起来的. \(^{1,2}\) 虽然搜索、辨识、策略形成以及这些活动的整合等算法可以在传统数字计算机上通过仿真进行建立和测试，但尚无硬件支持学习机所需的复杂计算结构. 主要问题是设计一种活性存储元件，其存储值在长期内保持稳定，能够以小的增量进行变化，且其输出可以作为``权重''来乘其他变量. 由于任何实用系统都需要大量此类元件，因此还要求它们体积小、成本低. 传统的模拟积分器和乘法器无法满足稳定性和低成本的要求，而非传统元件如电化学存储器和磁通变换器则不可靠，或需要复杂的外部电路才能使用. \(^{3}\) 半导体集成电路在速度、稳定性、尺寸和成本方面具有优势，因此决定设计一种基于标准门电路和触发器的计算元件，以实现大规模集成.

二进制双向计数器具有增量存储器的特性，但如果增量过小且大小可变，则需要很多位. 然而，若增量过程为随机过程，则可在存储计数中实现``分数增量''. 例如，若增量只需最低有效位值的十分之一，则计数器可以十分之一的概率增加一——这就是随机计算的基本原理：用事件发生的概率来表示模拟量. 这一原理最初在 ADDIE（平滑存储器件）中得到体现，它是一种具有随机输入和输出的平滑与存储器件，用于实现 STeLLA 学习方案，\(^{1}\) 后来扩展形成了一系列能够执行模拟计算机所有正常功能的计算元件家族——这被称为随机计算机. \(^{4}\)

由于本文篇幅所限，只能简要描述一些基本的随机计算元件和配置. 文献 4 讨论了随机计算的理论基础，并描述了一些先前在数据处理中使用随机变量的应用，而文献 5 则描述了随机计算在过程辨识中的应用，包括梯度技术、Bayes 估计与预测以及 Markov 建模.
